% !TeX root = ../../../graduate.tex
%% 第六章

\chapter{总结与展望}
\chaptermark{总结与展望}
\label{chap:conclusion}

本文研究外部高性能微内核如何进入 MLIR 调度和自动调优系统。链接外部符号只是执行链的最后一步。在此之前，编译器需要判断调度是否满足微内核的真实调用条件，并在条件不足时选择适配、回退或拒绝。候选能够正确调用微内核之后，还要单独回答它是否快于完整库，以及来源平台上的排序能否用于新目标。围绕这些问题，本文以微内核可接纳性契约为核心，完成知识建模、候选生成、真实执行、跨目标验证和搜索经济性分析。

\section{研究工作总结}

研究首先固定 BLIS 和 OpenBLAS 源码版本，从架构配置、微内核入口、packing 例程和调用代码中提取字段级事实。每条事实保留源码位置、证据等级、适用范围和核验状态，来源冲突按库与目标分别记录。决定调用能否成立的字段进入硬 Contract；缓存分块、参数偏好和候选排序进入可撤销的 Performance Profile。系统据此建立 BLIS Haswell \texttt{f32} $6\times16$、BLIS Haswell \texttt{f64} $6\times8$ 和 OpenBLAS RV64GCV \texttt{f32} $8\times8$ 三个可执行微内核包。

Contract 编译器根据任务上下文选择微内核包，随后传播 dtype、ISA、分块、packing、stride、tail 和 ABI 约束。非 reject 候选构成可执行集合 $\mathcal{Q}$，其中 direct/adapt 候选构成可调用子集 $\mathcal{C}$；fallback 保留正确的通用路径。候选记录还包含 Transform 配方、活跃参数、规范键、诊断和来源信息。路径相关的参数活性分析在调优前合并已知别名，避免把生成相同代码的参数组合重复交给搜索器。

调优器只消费有限的调度令牌，不修改 BaCO 的搜索核心。令牌被选中后，Transform Dialect 实例化调度，系统生成 A/B 面板、处理动态 M/N/K 边界、绑定外部符号，并继续降低到 LLVM IR 和目标二进制。BLIS 与 OpenBLAS 保留各自的 packing 格式和内部 ABI。候选评价同时检查数值结果、Transform/LLVM IR、链接符号、实际运行路径和边界计数，静态预测不能替代真实执行证据。

实验按 RQ1 至 RQ4 分别评价接纳与诊断、空间与搜索、迁移边界和部署经济性。静态分类、manifest 重放、冻结目标值上的离线搜索和物理机测量采用不同口径；完整 BLIS 或 OpenBLAS 始终作为部署性能边界。这样的证据安排使接口正确性、候选排序和搜索成本保持分离。

\section{主要结论}

\subsection{硬接纳边界可以成为独立的编译知识对象（RQ1）}

三个微内核包的独立变异实验包含 24 个危险单字段负例，全部被拒绝或安全降级，误接纳与误拒绝均为 0。正式执行中的路径预测也与运行观测一致。这些结果说明，微内核调用前提可以脱离库内分派代码，成为可审计、可诊断并由编译器消费的对象。

Contract 的能力仍有清楚边界。它检查声明字段是否匹配，不检查打包缓冲区的内容是否被破坏。实验中的两个内容故障由运行时数值参照检出。Contract、数值正确性和内存安全机制分别承担不同职责，不能相互替代。

\subsection{硬约束改善空间结构，软排序由目标环境决定（RQ2）}

硬 Contract 将 2,600 点原始域压缩到平均 104.4 个规范化唯一候选，并把 direct/adapt 密度提高到 1。它减少的是不满足外部接口或不改变代码的候选，作用可以通过诊断字段、规范化 IR 和真实路径解释。在 8 个工作负载各 28 点的冻结实测子池上，$D_{90}=0.1875$。同一 25 唯一查询预算下，BaCO 与均匀无放回随机搜索的中位 $T_{90}$ 为 7.5/7.0，达到率为 92.5\%/100\%；budget 5 与 10 的配对区间均跨 0，当前证据不支持 BaCO 优于随机搜索。该对照没有产生新的编译或硬件测量。packing 消融的负结果进一步表明，微内核接口义务会直接影响候选能否形成有效执行路径。

缓存分块和静态参数偏好没有表现出同样稳定的价值。原固定预算实验中 B3/B4 未优于 B2；严格同池实验中，目标数据能够改善 i7 的低预算搜索，当前 BLIS 软字段在 Data-only 之上没有统计支持的增量。性能排序应保留来源和目标条件，并允许在新测量到来后更新或撤销。

\subsection{可接纳性语义比性能排序更容易迁移（RQ3）}

在三个受测微内核包中，34 个 Schema 叶路径的交并比为 0.912，说明当前字段集合能够同时容纳这些对象。该比例只是受测对象的描述量，不是对未知后端字段重合度的总体估计。BLIS \texttt{f32}/\texttt{f64} 是同库、同 ISA、同后端的两个包，其中第二种数据类型复用了核心生成器；OpenBLAS RVV 是当前唯一独立库与独立 ISA 包，虽然完成了物理执行，候选生成和 packing 仍由包内专用代码承担。现有结果因此是分层迁移边界的存在性证据，不是通用迁移定律。

跨主机和跨 ISA 结果显示，合法执行可以保持，性能方向却会改变。i7 的校准收益没有在 i9 上复制，x86 来源的固定候选也不能支配 K230 工作负载。较稳定的迁移对象是明确的接纳边界和验证方法；候选优先级仍要依赖目标测量。

\subsection{能够搜索不等于值得搜索（RQ4）}

候选池内收敛只说明调优器能够找到该池中的较好配置。部署时还需比较完整库，并支付校准、编译、测量和调优成本。在 i9 和 K230 的预声明精确形状场景中，最佳生成路径相对完整库的性能比分别为 0.5448 和 0.6730，10 个工作负载中均为 0 个获胜。即使把搜索成本设为 0，生成路径也没有正向收益。

i7 上只有查看结果后筛出的 19 个有利形状在很长的合成调用规模上出现正净值。点估计盈亏平衡点约为 165,865 个等频合成周期，保守 LCB95 则需要约 741,587 个周期。这个事后场景只能说明成本量级，不能作为普遍部署收益。实际系统应在搜索前根据完整库基准、收益置信下界和预期调用量决定是否调优。

\section{局限性}

\subsection{知识抽取仍依赖固定版本和人工核验}

本文已经记录源码版本、字段来源、冲突与人工核验状态，但事实抽取尚未覆盖库内所有条件分派。BLIS 和 OpenBLAS 的新版本可能修改微内核、packing 或 ABI，使旧 Contract 失效。当前版本失效机制可以阻止继续使用过期对象，重新核验仍需要源码分析和人工判断。

三个微内核包不足以证明 Schema 对所有高性能库完备，0.912 也不能外推为新包的预期复用比例。BLIS \texttt{f32}/\texttt{f64} 属于同一后端；OpenBLAS RVV 提供了第二个库与第二种 ISA，却保留包内专用生成器。若要建立更强的后端通用性结论，还需接入具有不同调用组织和内存层次的第三种独立实现，并以前瞻冻结的缺失字段和接入修改记录检验 Schema，而不是只增加设备数量。

\subsection{系统覆盖的执行语义有限}

当前实现面向 rank-2、行主序、单线程 GEMM，支持受测的动态 M/N/K 与有限边界块。任意 stride、非连续布局、更多 $\alpha/\beta$ 模式、批量 GEMM、多线程和 NUMA 尚未形成统一 Contract 与适配器。完整 BLIS/OpenBLAS 还包含线程划分、宏观循环、微内核选择和更多边界优化，本文生成路径没有复现这些能力。

静态 checker 只处理声明元数据。别名关系、缓冲区生命周期、越界写入和内容损坏仍需更强的编译期分析、运行时保护或数值参照检查。当前两项故障注入证明数值门能够发现受测错误，没有覆盖所有内存安全问题。

\subsection{性能实验的主机与工作负载范围有限}

性能结论来自两台 x86 主机和一块 K230。固定预算区间重采样工作负载或随机种子，不包含机器层级的不确定性。i7 主机运行在 WSL2，尽管正式实验采用绑核和多进程重复，虚拟化环境仍会影响计时稳定性；未通过资格审计的数据已经排除，但无法把剩余结果视为原生 Linux 的完全替代。

Qwen 跟踪覆盖 Transformer block 内核心投影的 192 个规范调用，占全部 Linear FLOP 的 72.44\%。该过程没有执行完整模型，也没有包含 attention、激活、框架调度、内存分配或通信。按跟踪频次加权的结果只能说明受测 GEMM 调用的性能边界。

\subsection{校准和成本模型仍较粗}

目标数据校准使用 591 条测量，成本无法在 12 个 holdout 工作负载内摊销。当前日志能分解模型拟合、聚合测量和唯一编译成本，不能可靠分配每个形状的 $C_{s,j}$。价值模型因此适合判断明确的正负边界，不足以支持细粒度在线资源调度。

本文采用 LCB95 作为保守决策门，但未同时建模能耗、代码体积、内存占用和服务级延迟。调用频率来自冻结跟踪或合成场景，真实部署中的阶段变化、输入分布漂移和缓存驻留会改变盈亏平衡点。

\section{后续工作}

\subsection{把更多微内核包纳入共享生成接口}

近期最直接的工作是将 OpenBLAS RVV 的包内专用生成器收敛到共享核心接口。共享部分只消费 Contract 的语义字段和适配器能力，不加入按库名或微内核包标识的分支；RVV 面板、内部 ABI 和部署代码继续留在包内。完成后可比较修改前后的核心哈希，并结合独立包测试和物理机路径证据，检验后端特化是否真正减少。

随后可接入第三种独立实现，例如采用不同微内核生成或调用组织的 CPU 库。新增微内核包应先验证 Schema 的缺失字段，再讨论生成器复用，避免为了追求统一而丢失后端真实约束。

MLIR 接入也应沿第四章的分层边界演进：Contract Schema 与空间编译器保持独立，包内结构准备通过 registered pass 接口接入，专用 Transform operation 只承担无法由稳定上游原语清楚表达的精确重定向。后续版本以固定 LLVM revision、IR 回归和失败传播测试跟踪上游变化，不预设当前扩展一定适合直接上游化。

\subsection{自动化事实提取与版本失效检查}

现有来源记录已具备源码位置、版本和哈希，可以进一步构建增量核验工具。源码更新后，工具先定位受影响字段及其依赖 Contract，重新运行机器检查，再把存在歧义或冲突的项目交给人工审阅。对于由编译配置决定的 ISA、内核符号和分块大小，可从构建系统与目标对象中交叉提取，降低手工同步成本。

Contract Schema 还可扩展别名关系、缓冲区生命周期、对齐证明、线程约定和更完整的 $\alpha/\beta$ 语义。新增字段应保持可检查性；无法从源码或运行证据确认的内容继续进入 Unresolved，不填入推测值。

\subsection{降低目标校准和候选评价成本}

当前校准收益主要来自目标数据，瓶颈是获取这些数据的成本。后续可研究分阶段测量：先用低成本静态特征或短计时排除明显劣势候选，再对边界候选使用完整独立进程测量。已有编译产物可以按规范化 IR 和目标特征缓存，减少同义调度与相邻工作负载的重复成本。

搜索器还可直接优化净价值，而非只优化微内核运行时间。目标函数可以联合预测候选收益、编译成本、测量方差和剩余调用规模。当完整库的保守优势已经足够大时，系统应提前停止搜索；只有收益下界可能为正时才扩大预算。

\subsection{扩展布局、线程和算子范围}

GEMM 的下一步是补齐非连续 stride、转置、批量和多线程 Contract。多线程会引入线程分块、共享 packing、NUMA 和库内部并行冲突，不能只在现有单线程 Contract 上增加一个线程数参数。相关字段需要进入任务上下文、适配器能力和观测记录，并以完整多线程库为基准重新评价。

更广的算子扩展可以从具有明确微内核边界的批量 GEMM、卷积和 attention projection 开始。Contract 仍用于表达外部实现的接纳条件，算子融合和跨算子布局则需要增加区域级义务。现有共享激活 packing 实验没有超过完整库基线，说明新方向需要引入超出单次 packing 复用的机制，例如融合 epilogue、低精度布局或分阶段微内核组合。

\section{结语}

本文方法不替代成熟 GEMM 库，也不把固定参数视为可跨平台迁移的性能结论。它在编译器中显式保存微内核约束、适配能力和安全退出条件，使调优从可核验的接纳空间出发。候选通过正确性与路径检查后，仍须以目标测量确定排序，并与完整库按全成本比较；只有收益的保守下界能够覆盖这些成本时，搜索结果才具备部署意义。
